\subsection{Engineering application to steel surface defect clustering}
\label{subsec:neu_application}

To examine the engineering relevance of conservative graph calibration, we
considered semi-supervised clustering of hot-rolled steel-strip surface
defects using the public NEU-CLS database. The database contains six defect
categories with 300 grayscale images per category. Each image was resized to
$32\times32$, vectorized into a nonnegative column, and normalized to unit
Euclidean norm. No pretrained representation or external classifier was
used.

We evaluated label fractions of $5\%$, $10\%$, and $20\%$ per class. For
every fraction, 20 paired trials used identical labeled samples and initial
factors across GNMFLD, GOCNMF, and CGC-GOCNMF. Clustering labels were
obtained directly by the largest coordinate of the learned representation.
ACC, NMI, and ARI were computed only on unlabeled samples. We additionally
recorded the calibration coefficient, exact-fallback frequency, and aggregate
class-normalized confusion matrices.

The Stage-1 protocol fixes $p=5$ and $\alpha=10$ for GOCNMF and
CGC-GOCNMF, and uses $p=5$, $\alpha=10^{4}$, and $\beta=10$ for GNMFLD.
These are pre-specified pilot settings and are not claimed to be NEU-optimal.
No unlabeled ground-truth label is used in parameter selection.
